\chapter{Completeness of the secant dilates and the M\"obius problem}\label{chap:lv}\chaptermark{Completeness of the secant dilates}

Chapter~\ref{chap:rzeta} writes the exponential as a M\"obius weighted sum of hyperbolic secants (Proposition~\ref{rz:prop:lvj}) and records the question of Lai and Vere-Jones \cite{LVJ} on the completeness of the dilates $\sech(\lambda x)$ as open (Problem~\ref{rz:prob:nymanbeurling} asks for its relation to the Nyman Beurling criterion). This chapter answers the question in weighted $L^2$ spaces on the half line and says exactly where arithmetic enters. The answer has two halves. In one weight the integer dilates $\sech(nx)$ are complete unconditionally, by Carleman's formula, whatever the zeros of $L(s,\chi_4)$ do. In two weights at once completeness is equivalent to a zero free half plane of $L(s,\chi_4)$, with the generalised Riemann hypothesis at the threshold. The second half is a Nyman Beurling type criterion for the Dirichlet beta function, with an analogue for $\zeta$.

The chapter develops two source notes by Nick Polson, one on the secant dilates, Lerch's theorem and the M\"obius function (Sections~\ref{lv:sec:identity} to~\ref{lv:sec:thorin}) and one that works its four open problems (Sections~\ref{lv:sec:nec} to~\ref{lv:sec:open}), together with the computations done for this book (\texttt{code/lv\_sech\_tables.py}, \texttt{code/lv\_gram.py} and their outputs \texttt{data/lv\_sech\_tables.json}, \texttt{data/lv\_gram.json}, listed in \texttt{code/README.md}; the scripts of the sources were not used). Every number printed is one of these recomputations, and where the source differed the recomputed value is printed. The completeness question of Lai and Vere-Jones is posed in \cite{LVJ} as recorded in Chapter~\ref{chap:rzeta}; the paper itself was not consulted for this chapter, and no statement here is attributed to it beyond the question and the identity of Proposition~\ref{rz:prop:lvj}.

Notation. The book's $\Re$, $\Im$, $\sech$, $\R$, $\C$ and the cut plane $\cut$ are used. For $c>0$ the weighted space is $\mathcal W_c=L^2((0,\infty),x^{2c-1}dx)$; $\beta(s)=L(s,\chi_4)=1-3^{-s}+5^{-s}-\cdots$ is the Dirichlet beta function, a symbol local to this chapter (the real part of a zero is always written $\Re\rho$); $\mu$ is the M\"obius function, $M_\chi(x)=\sum_{n\le x}\mu(n)\chi_4(n)$ and $M(x)=\sum_{n\le x}\mu(n)$. The Dirichlet polynomial $\mathcal A_N(s)=\sum_{n\le N}\mu(n)\chi_4(n)n^{-s}$ is written with a script letter because $A_N$ is the coefficient of Theorem~\ref{rz:thm:sech}. The completed beta function is $\xi(s,\chi_4)$ of Section~\ref{sec:dbeta:beta}. The probability measures of one pair of zeros (Section~\ref{lv:sec:thorinform}) are written $Q_{a,\gamma}$. Status words follow the rest of the book: proved, computed, conditional, heuristic, open. Results that rest on the quasi Riemann hypothesis \cite{OAI26a,OAI26b,Liu26a} are conditional on it, and on one further point: the preprint is announced for Dirichlet $L$ functions, and whether its theorem covers the character modulo $4$ has not been checked for this book.

The results are these.
\begin{enumerate}
\item The integer dilates $\{\sech(nx)\}$ are complete in $L^1$ with the weight of Lerch's theorem, and in every $\mathcal W_c$, $c>0$, unconditionally; more generally $\{\sech(\lambda x)\}_{\lambda\in\Lambda}$, $\Lambda\subset[1,\infty)$, is complete in $\mathcal W_c$ whenever $\sum_\Lambda1/\lambda=\infty$ (Theorems~\ref{lv:thm:L1} and~\ref{lv:thm:single}, proved). The proof of the second statement replaces Lerch by Carleman's formula and uses no arithmetic.
\item Completeness in two weights $\mathcal W_a\cap\mathcal W_b$ for all $\theta<a<b<1$ is equivalent to $L(s,\chi_4)\ne0$ for $\Re s>\theta$, and to convergence of the M\"obius secant series in every $\mathcal W_c$, $\theta<c<1$ (Theorem~\ref{lv:thm:two}, proved); completeness in two weights straddling a strip forces that strip to be zero free (Proposition~\ref{lv:prop:nec}, proved). The zeta kernel $1/(e^x-1)-1/x$ has the same two statements, with $\zeta$ and RH (Theorem~\ref{lv:thm:zeta}, proved). The consequences of the quasi Riemann hypothesis are in Section~\ref{lv:sec:qrh} (conditional).
\item Under a polynomial bound for $1/\beta$ on the line $\Re s=c$, which holds for $c\ge1$, conditionally on a zero free half plane for $c$ to its right, and for $c>0$, $c\ne\frac12$ under GRH, the condition $\sum_\Lambda1/\lambda=\infty$ is also necessary (Theorem~\ref{lv:thm:nec}, proved); the line $c=\frac12$ stays open. The $L^1$ criterion of Salem for the integer dilates, and its beta analogue, are Theorem~\ref{lv:thm:salem} (proved).
\item The optimal distance of $e^{-x}$ from the span of $\sech(nx)$, $n\le N$, is given by a Gram system (computed, Table~\ref{lv:tab:dN}); the optimal coefficients are huge, and the explanation by the frequencies $\log n$ is heuristic (Section~\ref{lv:sec:gram}). The residual identity of Proposition~\ref{lv:prop:segment}, with its factor $2$, is proved.
\item For the Meixner family of order $t$ the Mellin transform is $M_t$ of Theorem~\ref{lv:thm:meixner}, single weight completeness holds throughout, and for $t=2$ the M\"obius route is unconditional (Theorem~\ref{lv:thm:t2}, proved).
\item For $\xi(s,\chi_4)$ the equivalence of the zero free half plane with a generalised gamma convolution form of $\xi(\theta,\chi_4)/\xi(\theta+\sqrt s,\chi_4)$ and the explicit Thorin measure are Theorem~\ref{lv:thm:thorin} (proved; for $\xi$ this is Theorem~\ref{thm:poisson}); the link with the two weight criterion is new (Section~\ref{lv:sec:thorin}).
\end{enumerate}
The open problems close the chapter.
